\begin{abstract}
Many waste-sorting papers classify a single object on a plain background, and papers that tune such models with a metaheuristic rarely compare against random search or report confidence intervals. We test whether particle swarm optimization (PSO) gives a better waste detector than random search or the library defaults when every method gets the same budget. Six training hyperparameters of YOLOv8n are tuned. Random search tries 32 settings and PSO runs 8 particles for 4 rounds, so each search costs 32 short proxy trainings. Before searching, we checked that the proxy ranks settings in the same order as longer training (Spearman correlation 0.976). The best setting of each method was then trained for 100 epochs and scored once on a held-out test set, with paired-bootstrap 95\% confidence intervals. We ran the experiment twice: first on studio images (\genone), then with real-world photographs added and a separate Indian street test set (\gentwo). On studio data the defaults, random search and PSO reach a mean average precision at IoU 0.5 (\mapfifty) of 0.676, 0.672 and 0.668, a statistical tie. In \gentwo both searches beat the defaults on the mixed test set by a small but significant margin (0.017 and 0.020) and do not differ from each other. On real-world photographs the three score 0.298 to 0.310 with no significant difference, and on Indian photographs only 0.023 to 0.027. Real-world training data mattered far more than tuning. On the same test images it raised real-world \mapfifty from 0.038--0.096 to 0.338--0.353, while studio \mapfifty fell by 0.074 to 0.105.
\end{abstract}

\begin{IEEEkeywords}
Waste detection, particle swarm optimization, hyperparameter optimization, YOLOv8, random search, domain shift.
\end{IEEEkeywords}
